\chapter{2022年度}
\label{ch:2022}

\section{第1期・専門基礎科目（3題必修）}

\begin{problem}{第1問｜実行列}{2022-1-1}
$n$ 次の単位行列を $E_n$ で表す。このとき，以下の問いに答えよ。
\begin{enumerate}
  \item[(1)] $A^2=-E_n$ となる $n$ 次の実正方行列 $A$ が存在するならば，$n$ は偶数であることを示せ。また，$A^2=-E_2$ となる2次の実正方行列 $A$ の例を挙げよ。
  \item[(2)] ${}^{\mathsf T}AA=-E_n$ となる $n$ 次の実正方行列 $A$ は存在しないことを示せ。
  \item[(3)] $A$ は零行列ではない $n$ 次の実正方行列で ${}^{\mathsf T}A=-A$ を満たすとする。$P^{-1}AP$ が対角行列となるような $n$ 次の実正則行列 $P$ は存在しないことを示せ。
\end{enumerate}
\end{problem}
\begin{memo*}
\textbf{解答の見通し。}
行列の等式は成分が多く、そのままでは比較しにくい。そこで、行列式や
$x^{\mathsf T}(\cdot)x$ を使って、行列の等式を一つの実数の等式へ変える。
(1) では $|BC|=|B|\,|C|$ により $A^2=-E_n$ を
$|A|^2=(-1)^n$ へ変え、左右の符号を比べる。
(2) では $x^{\mathsf T}A^{\mathsf T}Ax=\norm{Ax}^2\ge0$ なのに、
右辺は $-\norm{x}^2<0$ になることを使う。
(3) の $A^{\mathsf T}=-A$ は交代行列の条件であり、まず全ての実固有値が0だと示す。
もし実数上で対角化できれば、対角成分は全て0となって $A=O$ になり、仮定と矛盾する。
\end{memo*}
\begin{proof}[解答]
\begin{enumerate}
  \item[(1)] 行列式を取ると
  \[
  (\det A)^2=\det(-E_n)=(-1)^n.
  \]
  左辺は正だから $n$ は偶数である。また
  \[
  A=\begin{pmatrix}0&-1\\1&0\end{pmatrix}
  \]
  とすれば $A^2=-E_2$ である。

  \item[(2)] $0\ne x\in\R^n$ に対して
  \[
  x^{\mathsf T}A^{\mathsf T}Ax=\|Ax\|^2\ge0,\qquad
  x^{\mathsf T}(-E_n)x=-\|x\|^2<0.
  \]
  従って $A^{\mathsf T}A=-E_n$ は不可能である。

  \item[(3)] $Av=\lambda v$ $(0\ne v\in\R^n,\ \lambda\in\R)$ なら
  \[
  \lambda\|v\|^2=v^{\mathsf T}Av
  =(v^{\mathsf T}Av)^{\mathsf T}
  =v^{\mathsf T}A^{\mathsf T}v=-v^{\mathsf T}Av,
  \]
  よって $\lambda=0$ である。実数上で対角化できれば対角行列は零行列となり、
  $A=O$ となるので仮定に反する。
\end{enumerate}
\end{proof}

\begin{problem}{第2問｜Cauchy の関数方程式}{2022-1-2}
$\R$ 上で定義された実数値関数 $f(x)$ が，任意の実数 $x,y$ に対して
\[
f(x+y)=f(x)+f(y)
\]
を満たすとする。このとき，以下の問いに答えよ。
\begin{enumerate}
  \item[(1)] $f(0)=0$, $f(-x)=-f(x)$ が成り立つことを示せ。
  \item[(2)] 任意の有理数 $r\in\Q$ に対して，$f(rx)=rf(x)$ が成り立つことを示せ。
  \item[(3)] $f(x)$ が $\R$ 上で連続とすると，$f(x)$ はある定数 $c\in\R$ を用いて，$f(x)=cx$ と書けることを示せ。
\end{enumerate}
\end{problem}
\begin{memo*}
\textbf{解答の見通し。}
条件 $f(x+y)=f(x)+f(y)$ は、入力の足し算が出力の足し算へ移るという意味である。
まず式へ $x=y=0$ を入れて $f(0)$ を決め、次に $y=-x$ を入れて負号の扱いを決める。
同じ $x$ を繰り返し足せば整数倍、$x=n(x/n)$ と書けば $1/n$ 倍、両者を組み合わせれば
有理数倍に対する式が得られる。最後に、任意の実数 $x$ へ近づく有理数列を取り、
連続性によって有理数上の公式 $f(r)=f(1)r$ を実数全体へ移す。
\end{memo*}
\begin{proof}[解答]
\begin{enumerate}
  \item[(1)]
  \[
  f(0)=f(0)+f(0)\quad\Longrightarrow\quad f(0)=0,
  \]
  \[
  0=f(x-x)=f(x)+f(-x)\quad\Longrightarrow\quad f(-x)=-f(x).
  \]

  \item[(2)] 加法性と (1) より $f(mx)=mf(x)$ $(m\in\Z)$。また
  \[
  f(x)=f\left(n\frac{x}{n}\right)=nf\left(\frac{x}{n}\right)
  \]
  だから、$r=m/n$ $(m\in\Z,\ n\in\N)$ に対して
  \[
  f(rx)=f\left(m\frac{x}{n}\right)=\frac{m}{n}f(x)=rf(x).
  \]

  \item[(3)] $c=f(1)$ とおくと (2) より $f(r)=cr$ $(r\in\Q)$。
  $r_n\in\Q$, $r_n\to x$ を取れば、連続性から
  \[
  f(x)=\lim_{n\to\infty}f(r_n)
  =c\lim_{n\to\infty}r_n=cx.
  \]
\end{enumerate}
\end{proof}

\begin{problem}{第3問｜連続写像の反例}{2022-1-3}
位相空間 $X$ から位相空間 $Y$ への連続写像 $f:X\to Y$ に関する次の各命題はどれも成り立たない。反例を挙げよ。
\begin{enumerate}
  \item[(1)] $X$ の開集合 $U$ に対して像 $f(U)$ は $Y$ の開集合である。
  \item[(2)] $Y$ のコンパクト部分集合 $K$ に対して逆像 $f^{-1}(K)$ は $X$ のコンパクト部分集合である。
  \item[(3)] $f$ が全射のとき，$Y$ の稠密な部分集合 $V$ に対して逆像 $f^{-1}(V)$ は $X$ の稠密な部分集合である。
\end{enumerate}
\end{problem}
\begin{memo*}
\textbf{解答の見通し。}
反例とは、仮定を全て満たすのに結論だけが失敗する具体例である。
連続性が直接保証するのは「開集合の\emph{逆像}が開」であり、開集合の像や
コンパクト集合の逆像ではない。(1), (2) は全ての点を0へ送る定数写像で、この違いを示す。
(3) では全射性も必要なので、実数直線はそのまま恒等的に写し、別に加えた孤立点 $p$ だけを0へ送る。
孤立点とは一元集合 $\{p\}$ 自身が開になる点であり、これが逆像の稠密性を壊す。
\end{memo*}
\begin{proof}[解答]
\begin{enumerate}
  \item[(1)] 定数写像 $f:\R\to\R$, $f(x)=0$ は連続である。
  開集合 $U=\R$ に対して $f(U)=\{0\}$ は開でない。

  \item[(2)] 同じ写像で $K=\{0\}$ はコンパクトだが
  $f^{-1}(K)=\R$ はコンパクトでない。

  \item[(3)] $X=\R\sqcup\{p\}$ に直和位相を入れ、$Y=\R$ とする。
  \[
  f(x)=x\ (x\in\R),\qquad f(p)=0
  \]
  とおけば $f$ は連続かつ全射である。$V=\R\setminus\{0\}$ は $Y$ で稠密だが
  \[
  f^{-1}(V)=\R\setminus\{0\}
  \]
  は $X$ の非空開集合 $\{p\}$ と交わらないので稠密でない。
\end{enumerate}
\end{proof}

\section{第1期・専門科目（4題から2題選択）}

\begin{problem}{第1問｜有限群}{2022-s-1}
以下の問いに答えよ。
\begin{enumerate}
  \item[(1)] $\langle x\rangle$ を位数5の群とする。$\langle x\rangle$ の自己同型群（$\langle x\rangle$ から $\langle x\rangle$ への全単射準同型写像全体の集合）の位数を求めよ。答だけでなく，どのように求めるかも説明せよ。
  \item[(2)] $\{1,2,3,4,5\}$ に作用する5次対称群を $G$ とする。$P=\langle(1\ 2\ 3\ 4\ 5)\rangle$ とおく。$P$ の $G$ における正規化群 $N_G(P)=\{g\in G\mid gPg^{-1}=P\}$ のシロー2部分群を1つ求めよ。
  \item[(3)] 5次対称群 $G$ のシロー5部分群全体の集合を $\Omega$ とし，$\Omega$ の置換全体の集合を $\operatorname{Sym}(\Omega)$ とおく。$g\in G$ に対して
  \[
  \varphi(g):\Omega\longrightarrow\Omega,\qquad
  P\longmapsto\varphi(g)(P)=gPg^{-1}
  \]
  とおくと，$\varphi(g)\in\operatorname{Sym}(\Omega)$ である。これにより写像 $\varphi:G\to\operatorname{Sym}(\Omega)$ が定義できる。このとき，$\varphi$ は準同型写像であることを示せ。
  \item[(4)] 6次対称群には指数6の部分群で互いに共役ではないものが存在することを示せ。
\end{enumerate}
\end{problem}
\begin{memo*}
\textbf{解答の見通し。}
位数5の部分群を数えると、$S_5$ が6個の Sylow $5$ 部分群へ作用することが分かる。
この共役作用から通常とは異なる $S_5\hookrightarrow S_6$ を作り、
最後に「6点へ推移的に作用するか、1点を固定するか」という共役で変わらない性質で
二つの指数6部分群を区別する。
\end{memo*}
\begin{proof}[解答]
\begin{enumerate}
  \item[(1)] 位数5の群 $\langle x\rangle$ は巡回群である。自己同型 $\varphi$ は生成元 $x$ の像
  $\varphi(x)$ によって一意に決まり、また自己同型は生成元を生成元へ写す。
  $\langle x\rangle$ の生成元は
  \[
  x,x^2,x^3,x^4
  \]
  の4個である。逆に $k=1,2,3,4$ に対して $x\mapsto x^k$ は自己同型を定めるので、
  自己同型は生成元の像で尽くされ
  \[
  |\operatorname{Aut}(\langle x\rangle)|=4.
  \]

  \item[(2)] $S_5$ の5-cycle の個数は、最初の文字を1に固定して残りを並べることで
  \[
  (5-1)!=24
  \]
  個である。一つの位数5の部分群には単位元でない5-cycle が4個あり、異なる位数5の部分群は
  単位元以外では交わらない。従って Sylow $5$ 部分群の個数は
  \[
  |\Omega|=\frac{24}{4}=6.
  \]
  共役作用における $P$ の安定化群は $N_G(P)$ であり、Sylow の定理により作用は推移的だから、
  軌道・安定化群定理から
  \[
  |N_G(P)|=\frac{|S_5|}{|\Omega|}
  =\frac{120}{6}=20=2^2\cdot5.
  \]

  ここで
  \[
  a=(1\ 2\ 3\ 4\ 5),
  \qquad b=(2\ 3\ 5\ 4)
  \]
  とおく。$b$ は4-cycle なので位数は4である。また巡回置換の共役公式から
  \[
  bab^{-1}
  =(b(1)\ b(2)\ b(3)\ b(4)\ b(5))
  =(1\ 3\ 5\ 2\ 4)=a^2.
  \]
  従って $bPb^{-1}=P$、すなわち $b\in N_G(P)$ である。
  よって $\langle b\rangle$ は $N_G(P)$ の位数4の部分群であり、
  $|N_G(P)|$ に現れる2の最高冪が $2^2=4$ だから Sylow $2$ 部分群である。

  \item[(3)] $g\in G$ と $Q\in\Omega$ に対して
  \[
  \varphi(g)(Q)=gQg^{-1}
  \]
  と定める。$Q$ が位数5なら $gQg^{-1}$ も位数5なので、$\varphi(g)$ は $\Omega$ の置換を定める。
  $g,h\in G$ に対して
  \begin{align*}
  \varphi(gh)(Q)
  &=(gh)Q(gh)^{-1}\\
  &=g(hQh^{-1})g^{-1}\\
  &=\varphi(g)(\varphi(h)(Q)).
  \end{align*}
  従って $\varphi(gh)=\varphi(g)\varphi(h)$ であり、
  $\varphi:G\to\operatorname{Sym}(\Omega)$ は群準同型である。

  \item[(4)] (2) より $|\Omega|=6$ なので、$\Omega$ と $\{1,\ldots,6\}$ を同一視すれば
  \[
  \varphi:S_5\longrightarrow S_6
  \]
  とみなせる。$K=\ker\varphi$ とおくと $K\triangleleft S_5$ である。
  $K$ の元は全ての $Q\in\Omega$ を固定し、特に $P$ を固定するから
  \[
  K\subset N_G(P).
  \]
  $S_5$ の正規部分群は $\{1\},A_5,S_5$ の三つである。
  一方 $|N_G(P)|=20$ なので、位数60の $A_5$ も位数120の $S_5$ も
  $N_G(P)$ に含まれない。従って $K=\{1\}$ であり、$\varphi$ は単射である。

  よって
  \[
  H_1:=\varphi(S_5)\subset S_6,
  \qquad |H_1|=120,
  \qquad [S_6:H_1]=\frac{720}{120}=6.
  \]
  Sylow 部分群は互いに共役なので、$H_1$ の6点集合 $\Omega$ への作用は推移的である。

  一方、通常の点固定部分群
  \[
  H_2:=\{\tau\in S_6\mid\tau(6)=6\}
  \]
  は $S_5$ と同型で、$|H_2|=120$、従って $[S_6:H_2]=6$ である。
  $H_2$ は点6を固定するが、$H_1$ は6点に推移的に作用するため固定点を持たない。
  もし $H_1=sH_2s^{-1}$ となる $s\in S_6$ が存在すれば、$H_1$ は点 $s(6)$ を固定することになり矛盾する。
  従って $H_1,H_2$ は指数6で互いに共役でない二つの部分群である。
\end{enumerate}
\end{proof}

\begin{problem}{第2問｜球面の面積形式}{2022-s-2}
$\R^3$ 内の球面 $S^2:=\{(x,y,z)\in\R^3\mid x^2+y^2+z^2=1\}$ を考える。点 $p\in S^2$ に対して，
\[
\omega_p(u,v):=\inner p{u\times v}
\]
と定める。ただし，$u,v$ は点 $p$ における接ベクトルであり，$u=(u_1,u_2,u_3)$, $v=(v_1,v_2,v_3)$ に対して $u\times v$ は
\[
u\times v=(u_2v_3-u_3v_2,\ u_3v_1-u_1v_3,\ u_1v_2-u_2v_1)
\]
で定義される外積である。また，$\inner{-}{-}$ は $\R^3$ の標準内積である。このとき，以下の問いに答えよ。
\begin{enumerate}
  \item[(1)] $\omega_p(u,v)=-\omega_p(v,u)$ が成り立つことおよび $\omega_p(u,v)=0$ が全ての接ベクトル $v$ で成り立つならば $u=0$ であることを示せ。
  \item[(2)] $S^2$ の開集合 $U$ を
  \[
  U=S^2\setminus\{(x,y,z)\in S^2\mid y=0,\ x\ge0\}
  \]
  で定める。$0<\theta<2\pi$, $-1<h<1$ に対し，
  \[
  (r(\theta,h)\cos\theta,\ r(\theta,h)\sin\theta,\ h)\in U
  \]
  となる連続関数 $r(\theta,h)$ を求めよ。
  \item[(3)] (2) で求めた $r(\theta,h)$ を用いて，写像 $T:(0,2\pi)\times(-1,1)\to U$ を
  \[
  T(\theta,h)=(r(\theta,h)\cos\theta,\ r(\theta,h)\sin\theta,\ h)
  \]
  で定める。$p=T(\theta,h)\in U$ に対して
  \[
  \omega_p\left(\frac{\partial T}{\partial\theta},\frac{\partial T}{\partial h}\right)
  \]
  を求めよ。
\end{enumerate}
\end{problem}
\begin{memo*}
\textbf{解答の見通し。}
(1) の反対称性は外積の順序交換から従う。非退化性では
$v=p\times u$ と選ぶと、三重積が $\norm{u}^2p$ になって $u$ の長さを取り出せる。
(2), (3) では高さ $h$ の緯線の半径を球面の方程式から決め、
座標ベクトルの外積が外向き法線 $T$ に一致することを確かめる。
\end{memo*}
\begin{proof}[解答]
\begin{enumerate}
  \item[(1)] 外積の反対称性 $v\times u=-(u\times v)$ より
  \[
  \omega_p(v,u)=\inner p{v\times u}
  =-\inner p{u\times v}=-\omega_p(u,v).
  \]
  次に $\omega_p(u,v)=0$ が全ての $v\in T_pS^2$ に対して成り立つと仮定する。
  $u\in T_pS^2$ だから $\inner pu=0$ であり、$v=p\times u$ も $p$ と直交するので
  $v\in T_pS^2$ である。ベクトル三重積の公式を用いると
  \[
  u\times(p\times u)
  =p\inner uu-u\inner up
  =\norm u^2p.
  \]
  従って
  \[
  0=\omega_p(u,p\times u)
  =\inner p{\norm u^2p}
  =\norm u^2\norm p^2=\norm u^2.
  \]
  よって $u=0$ である。

  \item[(2)] 点 $(r\cos\theta,r\sin\theta,h)$ が $S^2$ 上にある条件は
  \[
  r^2\cos^2\theta+r^2\sin^2\theta+h^2=1,
  \]
  すなわち $r^2=1-h^2$ である。$-1<h<1$ では右辺は正であり、
  $r$ を緯線の半径として正に取るので
  \[
  r(\theta,h)=\sqrt{1-h^2}.
  \]
  $r$ は $\theta$ に依存しない。また $0<\theta<2\pi$ としたことで、除いた半大円上の点は現れない。

  \item[(3)] 以下 $T_\theta=\partial T/\partial\theta$, $T_h=\partial T/\partial h$ と略記する。
  $r=\sqrt{1-h^2}$ だから $\partial r/\partial h=-h/r$ であり
  \[
  T_\theta=(-r\sin\theta,r\cos\theta,0),\qquad
  T_h=\left(-\frac h r\cos\theta,-\frac h r\sin\theta,1\right).
  \]
  外積を成分ごとに計算すると
  \begin{align*}
  T_\theta\times T_h
  &=\left(r\cos\theta,r\sin\theta,
    h(\sin^2\theta+\cos^2\theta)\right)\\
  &=(r\cos\theta,r\sin\theta,h)=T(\theta,h).
  \end{align*}
  従って $\omega_T(T_\theta,T_h)=\inner{T}{T}=1$。
\end{enumerate}
\end{proof}

\begin{problem}{第3問｜微分方程式}{2022-s-3}
$f(x)$ を $\R$ 上の有界かつ連続な関数とし，微分方程式
\[
\frac{\d^2y}{\d x^2}-y=f(x)
\tag{$*$}
\]
を考える。このとき，以下の問いに答えよ。
\begin{enumerate}
  \item[(1)] $f(x)=0$ $(x\in\R)$ とするとき，微分方程式 ($*$) の一般解を求めよ。
  \item[(2)] $a\in\R$ とする。このとき，$e^{-x}f(x)$ は区間 $[a,\infty)$ において広義積分可能であることを示せ。
  \item[(3)] $\R$ 上の関数 $g(x)$ を
  \[
  g(x)=-\frac{e^x}2\int_x^\infty e^{-t}f(t)\,\d t
  -\frac{e^{-x}}2\int_{-\infty}^xe^tf(t)\,\d t
  \qquad(x\in\R)
  \]
  によって定義する。このとき，$g(x)$ は微分方程式 ($*$) の解であることを示せ。
  \item[(4)] $\displaystyle\lim_{x\to\infty}f(x)=1$ を満たすとき，(3) で定義した関数 $g(x)$ は $\displaystyle\lim_{x\to\infty}g(x)=-1$ を満たすことを示せ。
  \item[(5)] 
  \[
  f(x)=\begin{cases}1-e^{-x}&(x\ge0)\\0&(x<0)\end{cases}
  \]
  のとき $\displaystyle\lim_{x\to\infty}y(x)=-1$ を満たす微分方程式 ($*$) の解を全て求めよ。
\end{enumerate}
\end{problem}
\begin{memo*}
\textbf{解答の見通し。}
最初に同次方程式の二つの基本解 $e^x,e^{-x}$ を確認する。
(3) の積分表示は二つの積分を別々に微分すると交差項が消えるよう設計されている。
(4) は $f=1+(f-1)$ と分けて優収束定理を使い、
(5) は $x<0$ と $x\ge0$ で解いた後、$x=0$ で $y,y'$ を接続する。
\end{memo*}
\begin{proof}[解答]
\begin{enumerate}
  \item[(1)] 特性方程式は
  \[
  \lambda^2-1=(\lambda-1)(\lambda+1)=0
  \]
  であり、相異なる特性根は $1,-1$ である。従って一般解は
  \[
  y(x)=C_1e^x+C_2e^{-x}\qquad(C_1,C_2\in\R)
  \]
  である。
  \item[(2)] $f$ は有界だから、ある $M\ge0$ が存在して
  $|f(x)|\le M$ が全ての $x$ で成り立つ。任意の $a\in\R$ に対して
  \[
  \int_a^\infty|e^{-x}f(x)|\,\d x
  \le M\int_a^\infty e^{-x}\,\d x
  =Me^{-a}<\infty.
  \]
  よって $e^{-x}f(x)$ は $[a,\infty)$ 上で絶対可積分であり、広義積分可能である。
  \item[(3)]
  \[
  A(x)=\int_x^\infty e^{-t}f(t)\,\d t,\qquad
  B(x)=\int_{-\infty}^x e^tf(t)\,\d t
  \]
  とおく。微積分学の基本定理より
  \[
  A'(x)=-e^{-x}f(x),\qquad B'(x)=e^xf(x).
  \]
  $g=-(e^xA+e^{-x}B)/2$ を積の微分法で微分すると
  \begin{align*}
  g'
  &=-\frac12\{e^xA+e^xA'-e^{-x}B+e^{-x}B'\}\\
  &=-\frac12e^xA+\frac12e^{-x}B,
  \end{align*}
  ここでは $e^xA'=-f$ と $e^{-x}B'=f$ が打ち消し合った。もう一度微分して
  \begin{align*}
  g''
  &=-\frac12(e^xA+e^xA')
    +\frac12(-e^{-x}B+e^{-x}B')\\
  &=-\frac12(e^xA+e^{-x}B)+f(x)\\
  &=g(x)+f(x).
  \end{align*}
  従って $g''-g=f$ となり、$g$ は微分方程式 ($*$) の解である。
  \item[(4)] $h=f-1$ とおくと $h(x)\to0$ で、$h$ は有界である。
  定数関数 $1$ を (3) の式へ代入して得られる解は
  \[
  -\frac{e^x}{2}\int_x^\infty e^{-t}\,\d t
  -\frac{e^{-x}}2\int_{-\infty}^x e^t\,\d t=-1.
  \]
  一方、$h$ が作る第1項は変数変換 $t=x+s$ により
  \[
  e^x\int_x^\infty e^{-t}h(t)\,\d t
  =\int_0^\infty e^{-s}h(x+s)\,\d s\longrightarrow0,
  \]
  第2項も $t=x-s$ により
  \[
  e^{-x}\int_{-\infty}^x e^th(t)\,\d t
  =\int_0^\infty e^{-s}h(x-s)\,\d s\longrightarrow0.
  \]
  いずれも有界性から $Ce^{-s}$ を優関数にして優収束定理を使った。
  従って $g(x)\to-1$。
  \item[(5)] $x\ge0$ では、定数項 $1$ の特殊解は $-1$。
  また $-e^{-x}$ は同次解 $e^{-x}$ と共振するので
  $y_p=Cxe^{-x}$ とおくと
  \[
  y_p''-y_p=-2Ce^{-x}.
  \]
  よって $C=1/2$。$x\to\infty$ で $-1$ に収束するには $e^x$ の項があってはならないから
  \[
  y(x)=-1+\left(D+\frac x2\right)e^{-x}\qquad(x\ge0).
  \]
  $x<0$ では $f=0$ なので $y=Be^x+Ce^{-x}$。
  $x=0$ で $y,y'$ を一致させると
  \[
  B+C=-1+D,\qquad B-C=\frac12-D,
  \]
  従って $B=-1/4$, $C=D-3/4$。以上より、任意の $D\in\R$ に対し
  \[
  y(x)=\begin{cases}
  -\frac14e^x+(D-\frac34)e^{-x},&x<0,\\
  -1+(D+\frac x2)e^{-x},&x\ge0
  \end{cases}
  \]
  であり、これが全てである。
\end{enumerate}
\end{proof}

\begin{problem}{第4問｜Dirichlet 積分}{2022-s-4}
以下の問いに答えよ。
\begin{enumerate}
  \item[(1)] $t$ を正の実数とする。$\dfrac{\sin x}x$ は開区間 $(0,t)$ 上でルベーグ積分可能であることを示せ。
  \item[(2)] $t,y$ を正の実数とする。等式
  \[
  \int_0^t e^{-xy}\sin x\,\d x
  =\frac1{1+y^2}-\frac{e^{-ty}}{1+y^2}(y\sin t+\cos t)
  \]
  を示せ。
  \item[(3)] $t$ を正の実数とする。等式
  \[
  \int_0^t\frac{\sin x}{x}\,\d x
  =\int_0^\infty\frac{\d y}{1+y^2}
  -\int_0^\infty\frac{e^{-ty}}{1+y^2}(y\sin t+\cos t)\,\d y
  \]
  \item[(4)] $\displaystyle\lim_{t\to\infty}\int_0^t\frac{\sin x}{x}\,\d x$ を求めよ。
\end{enumerate}
\end{problem}
\begin{memo*}
\textbf{解答の見通し。}
$\sin x/x$ は原点の除去可能特異点を埋めれば連続になる。
次に $1/x=\int_0^\infty e^{-xy}\,\d y$ を使って二重積分へ移し、
Fubini の定理で積分順序を交換する。
最後は明示式の誤差項が $1/t+1/t^2$ 以下になることを示せばよい。
\end{memo*}
\begin{proof}[解答]
\begin{enumerate}
  \item[(1)] $\lim_{x\to0}\sin x/x=1$ なので、$x=0$ で値を1と定めれば $[0,t]$ 上連続になる。
  従ってコンパクト区間上でルベーグ積分可能である。
  \item[(2)]
  \[
  F(x)=\frac{e^{-xy}(-y\sin x-\cos x)}{1+y^2}
  \]
  とおく。積の微分法により
  \begin{align*}
  F'(x)
  &=\frac{e^{-xy}}{1+y^2}
  \{(-y)(-y\sin x-\cos x)-y\cos x+\sin x\}\\
  &=e^{-xy}\sin x.
  \end{align*}
  よって
  \begin{align*}
  \int_0^t e^{-xy}\sin x\,\d x
  &=F(t)-F(0)\\
  &=\frac1{1+y^2}
   -\frac{e^{-ty}}{1+y^2}(y\sin t+\cos t).
  \end{align*}
  \item[(3)] $x>0$ では
  \[
  \frac1x=\int_0^\infty e^{-xy}\,\d y.
  \]
  また
  \[
  \int_0^t\int_0^\infty|e^{-xy}\sin x|\,\d y\d x
  =\int_0^t\frac{|\sin x|}{x}\,\d x<\infty
  \]
  なので Fubini の定理で積分順序を交換できる。(2) を用いて
  \begin{align*}
  \int_0^t\frac{\sin x}{x}\,\d x
  &=\int_0^\infty\int_0^t e^{-xy}\sin x\,\d x\d y\\
  &=\int_0^\infty\frac{\d y}{1+y^2}
  -\int_0^\infty\frac{e^{-ty}}{1+y^2}(y\sin t+\cos t)\,\d y.
  \end{align*}
  \item[(4)] 第2積分の絶対値は
  \[
  \int_0^\infty e^{-ty}(y+1)\,\d y
  =\frac1{t^2}+\frac1t\longrightarrow0
  \]
  以下である。一方
  \[
  \int_0^\infty\frac{\d y}{1+y^2}
  =[\arctan y]_0^\infty=\frac\pi2.
  \]
  従って求める極限は $\pi/2$。
\end{enumerate}
\end{proof}

\section{第2期・専門基礎科目（3題必修）}

\begin{problem}{第1問｜線形空間}{2022-2-1}
以下，$U,V$ は $\R$ 上の線形空間とする。
\begin{enumerate}
  \item[(1)] $\R$ 上の線形空間 $V$ に対して，次の問いに答えよ。\\
  (i) 線形空間 $V$ の部分空間 $W_1,W_2$ に対して，$W_1\cap W_2$ も $V$ の部分空間となることを示せ。\\
  (ii) 線形空間 $V$ の部分空間 $W_1,W_2$ に対して，$W_1\cup W_2$ が $V$ の部分空間となるとき，$W_1\subset W_2$ または $W_2\subset W_1$ であることを示せ。
  \item[(2)] $U,V$ は $\R$ 上の線形空間で，$f:U\to V$ は線形写像であるとする。$U,V$ の零ベクトルをそれぞれ $0_U,0_V$ で表す。$\ker f=\{x\in U\mid f(x)=0_V\}$ は $U$ の部分空間であることを示せ。また，$f$ が単射であるための必要十分条件は，$\ker f=\{0_U\}$ であることを示せ。
  \item[(3)] 実数係数の3次正方行列全体を $V=M_3(\R)$ とし，$A\in V$ が交代行列であるとは ${}^{\mathsf T}A=-A$ を満たすことをいう。交代行列全体の部分集合を $W=\{A\in V\mid {}^{\mathsf T}A=-A\}$ とする。$W$ が $V$ の部分空間であることを示せ。また $W$ の1組の基底を与え，$W$ の次元を求めよ。
\end{enumerate}
\end{problem}
\begin{memo*}
\textbf{解答の見通し。}
部分空間であることは、零ベクトルを含み、集合の任意の元 $x,y$ と任意の $a,b\in\R$ に対して
$ax+by$ も集合内に残ることを確認すればよい。
合併の問題では、どちらにも含まれないベクトルを一つずつ取ると、その和が矛盾を生む。
核 $\ker f$ は0へ写るベクトル全体である。単射の条件 $f(x)=f(y)\Rightarrow x=y$ を、
線形性により $f(x-y)=0$ へ直すと核との関係が見える。
交代行列では $A^{\mathsf T}=-A$ を各成分の式 $a_{ji}=-a_{ij}$ に直し、
対角成分と下三角成分が上三角成分から決まることを読む。
\end{memo*}
\begin{proof}[解答]
\begin{enumerate}
  \item[(1)] (i) $0\in W_1\cap W_2$ であり、$x,y\in W_1\cap W_2$、
  $a,b\in\R$ なら $ax+by\in W_1\cap W_2$。従って部分空間である。

  (ii) どちらの包含も成り立たないとすると
  $u\in W_1\setminus W_2$, $v\in W_2\setminus W_1$ を取れる。
  $u+v\in W_1$ なら $v\in W_1$、$u+v\in W_2$ なら $u\in W_2$ となる。
  従って $u+v\notin W_1\cup W_2$ であり、合併が部分空間であることに反する。

  \item[(2)] $x,y\in\ker f$, $a,b\in\R$ なら
  \[
  f(ax+by)=af(x)+bf(y)=0
  \]
  であり、$0_U\in\ker f$ だから $\ker f$ は部分空間である。また
  \[
  f(x)=f(y)\iff f(x-y)=0\iff x-y\in\ker f.
  \]
  よって $f$ が単射であることと $\ker f=\{0_U\}$ は同値である。

  \item[(3)] $O\in W$ であり、$A,B\in W$, $a,b\in\R$ なら
  \[
  (aA+bB)^{\mathsf T}=-aA-bB,
  \]
  従って $W$ は部分空間である。任意の $A\in W$ は一意に
  \[
  A=\begin{pmatrix}0&a&b\\-a&0&c\\-b&-c&0\end{pmatrix}
  =a(E_{12}-E_{21})+b(E_{13}-E_{31})+c(E_{23}-E_{32})
  \]
  と書ける。従って
  \[
  E_{12}-E_{21},\quad E_{13}-E_{31},\quad E_{23}-E_{32}
  \]
  が基底であり、$\dim W=3$。
\end{enumerate}
\end{proof}

\begin{problem}{第2問｜微分積分}{2022-2-2}
\begin{enumerate}
  \item[(1)] $a<b$ とする。閉区間の直積集合 $[a,b]\times[a,b]$ 上の連続関数 $u(s,t)$ と $x\in(a,b)$ に対し
  \[
  f(x)=\int_a^x\d s\int_a^su(s,t)\,\d t
  \]
  とおく。$f(x)$ は開区間 $(a,b)$ 上微分可能で，その導関数は
  \[
  \int_a^xu(x,t)\,\d t
  \]
  で与えられることを示せ。
  \item[(2)] $x>0$ に対し
  \[
  D_x=\{(s,t)\mid0\le s\le x,\ s\le t\le x\}
  \]
  とおく。このとき関数
  \[
  g(x)=\iint_{D_x}s(1-2st)\,\d s\d t
  \]
  の $x>0$ の範囲における最大値を求めよ。
\end{enumerate}
\end{problem}
\begin{memo*}
\textbf{解答の見通し。}
(1) では内側の積分全体を $F(s)=\int_a^s u(s,t)\,\d t$ と名付ける。
すると $f(x)=\int_a^xF(s)\,\d s$ なので、$F$ が連続だと確認できれば
微積分学の基本定理から $f'(x)=F(x)$ となる。
$F$ では積分区間と被積分関数の両方が $s$ とともに変わるため、差を二つに分けて連続性を示す。
(2) の不等式 $0\le s\le t\le x$ は、外側で $t$ を $0$ から $x$ まで動かし、
各 $t$ を固定した内側で $s$ を $0$ から $t$ まで動かす、という積分順序を表す。
\end{memo*}
\begin{proof}[解答]
\begin{enumerate}
  \item[(1)] $F(s)=\int_a^s u(s,t)\,\d t$ とおく。
  $u$ の一様連続性と有界性から、$h\to0$ のとき
  \[
  F(s+h)-F(s)
  =\int_a^s\{u(s+h,t)-u(s,t)\}\,\d t
   +\int_s^{s+h}u(s+h,t)\,\d t\longrightarrow0.
  \]
  従って $F$ は連続であり、微積分学の基本定理から
  \[
  f'(x)=F(x)=\int_a^x u(x,t)\,\d t.
  \]

  \item[(2)] $D_x=\{(s,t)\mid0\le s\le t\le x\}$ より
  \begin{align*}
  g(x)
  &=\int_0^x\int_0^t(s-2s^2t)\,\d s\d t\\
  &=\int_0^x\left(\frac{t^2}{2}-\frac{2t^4}{3}\right)\d t
  =\frac{x^3}{6}-\frac{2x^5}{15}.
  \end{align*}
  \[
  g'(x)=x^2\left(\frac12-\frac{2x^2}{3}\right)
  \]
  だから、$g$ は $x=\sqrt3/2$ で最大となる。その値は
  \[
  g\left(\frac{\sqrt3}{2}\right)=\frac{\sqrt3}{40}.
  \]
\end{enumerate}
\end{proof}

\begin{problem}{第3問｜位相空間}{2022-2-3}
\begin{enumerate}
  \item[(1)] $X$ を位相空間とし，$A,B$ を $X$ の部分集合とする。このとき $\overline{A\cup B}=\overline A\cup\overline B$ を示せ。但し $\overline A,\overline B,\overline{A\cup B}$ はそれぞれ $A,B,A\cup B$ の閉包である。
  \item[(2)] ハウスドルフ空間 $Y$ のコンパクトな部分集合 $K$ は閉集合であることを示せ。
  \item[(3)] $S,T$ を位相空間とし，$F:S\to T$ を全射な連続写像とする。$S$ が連結ならば，$T$ も連結であることを示せ。
\end{enumerate}
\end{problem}
\begin{memo*}
\textbf{解答の見通し。}
(1) の閉包 $\overline A$ は「$A$ を含む最小の閉集合」である。
二つの集合が等しいことは、左が右に含まれることと、右が左に含まれることを別々に示せばよい。
(2) の Hausdorff 性は異なる二点を交わらない開近傍で分けられる性質、
コンパクト性は任意の開被覆から有限個を選べる性質である。
点 $x\notin K$ と各 $y\in K$ をまず個別に分け、有限個の近傍だけで $K$ 全体を覆う。
(3) の連結とは、空でない二つの開集合へ分割できないことである。
像が分割できると仮定し、連続写像の逆像を取ると定義域まで分割されることを示す。
\end{memo*}
\begin{proof}[解答]
\begin{enumerate}
  \item[(1)] $\overline A\cup\overline B$ は $A\cup B$ を含む閉集合だから
  \[
  \overline{A\cup B}\subset\overline A\cup\overline B.
  \]
  一方、$A,B\subset A\cup B$ より
  $\overline A,\overline B\subset\overline{A\cup B}$。従って等号が成り立つ。

  \item[(2)] $x\notin K$ とする。各 $y\in K$ に対し Hausdorff 性より
  \[
  x\in U_y,\quad y\in V_y,\quad U_y\cap V_y=\varnothing
  \]
  となる開集合を取る。$K$ のコンパクト性より
  $K\subset V_{y_1}\cup\cdots\cup V_{y_m}$ とできる。このとき
  \[
  U=U_{y_1}\cap\cdots\cap U_{y_m}
  \]
  は $x$ の開近傍で $U\cap K=\varnothing$。従って $X\setminus K$ は開であり、
  $K$ は閉である。

  \item[(3)] $T=U\cup V$ が非自明な開分離だとする。連続性と全射性より
  $F^{-1}(U),F^{-1}(V)$ は空でない $S$ の開集合であり、
  \[
  S=F^{-1}(U)\cup F^{-1}(V),\qquad
  F^{-1}(U)\cap F^{-1}(V)=\varnothing.
  \]
  これは $S$ の連結性に反する。従って $T$ は連結である。
\end{enumerate}
\end{proof}
